\chapter{Limitations, Reproducibility, and Future Work}

The thesis ends by separating current conclusions from remaining work. The
model is useful for mechanism separation and surrogate-ready data generation,
but device-specific lifetime prediction still requires stronger calibration,
ion-sensitive validation, and broader experimental constraints.

\section{Limitations}

The COMSOL model is a physics-based accelerated degradation study, not a final
lifetime prediction for a specific commercial device. The
starting electronic behavior must still be calibrated against the measured
fresh-device J--V range, and the degradation rates are accelerated for
computational practicality. The reported aging time should therefore be
interpreted as equivalent accelerated aging time, not as a direct field-time
prediction.

The second limitation is that different causes can look similar in a J--V
curve. Defect growth, leakage, contact resistance, surface recombination, and
reduced light-generated current can all lower PCE. Several combinations can
produce similar measured J--V degradation. The model needs hysteresis, impedance,
transient-current, optical, and aging data to separate these mechanisms with
confidence.

The third limitation is ion validation. The normalized ion variable
$u_{\mathrm{ion}}$ represents ion movement inside the model, but ion mobility,
concentration, boundary behavior, and the link between ion motion and defect
growth need
experimental support before they can be treated as device-specific
material parameters. DIT and capacitance-based measurements provide stronger
constraints than ordinary J--V curves, but the integrated current must be
baseline-corrected so that electronic plateaus and capacitive spikes are not
mistaken for mobile ionic charge.

The fourth limitation is implementation consistency with the laboratory stack.
The physical architecture is
glass/ITO/NiO$_x$/MeO-2PACz/perovskite/C$_{60}$/BCP/Ag/UV-resin/glass,
illuminated from the ITO side. All stack labels, domain selections,
illumination direction, coordinate conventions, and contact work functions must
remain consistent with that architecture before the model is called
device-specific.

The fifth limitation is diagnostic-voltage completeness. The light-temperature
stress-matrix J--V study is physically scaled and usable for PCE and $J_{sc}$
across all 360 curves, but its 1.25 V diagnostic voltage limit does not resolve
zero-current crossings for all mild or early-aging cases.
Quantitative $V_{oc}$ and FF claims therefore remain limited to the subset of
curves that cross zero current. A wider diagnostic voltage sweep is needed
before those quantities can be reported for every case.

\section{Reproducibility}

The analysis is organized so that the reported figures and tables can be
regenerated from the COMSOL-derived tables. The reproducibility package
separates simulation outputs, processed tables, plotting routines, thesis
source, and generated figures. Appendix~C gives the practical execution
details for rebuilding the results without interrupting the scientific
narrative in the main chapters.

\section{Data Availability}

The thesis source, processed analysis artifacts, generated figures and tables,
and public PDF are maintained in the supporting public repository at
\url{https://github.com/flabio121/thesis-public}. Large simulation
intermediates are excluded from the public copy when they are too large for
practical version control. Appendix~C and the reproducibility files document
how the public artifacts map to the underlying simulation studies.

\section{Future Work}

The next step is improved calibration. The pristine COMSOL
baseline should be fit to the same device architecture and measurement
conditions as the experimental J--V data. The fit should match $V_{oc}$,
$J_{sc}$, FF, and PCE before degradation parameters are interpreted.

The second step is implementation and validation of light-soaking and
self-healing terms. The proposed conditioning state can represent improved
contact extraction under illumination, while optional healing terms can
represent partial recovery during rest. These terms require controlled
light-soaking, dark-rest, and repeated-scan experiments before they should be
used in predictive simulations.

The third step is ion-sensitive validation. Electrical impedance spectroscopy
(EIS), transient photocurrent (TPC), capacitance, ion-current, and
scan-rate-dependent hysteresis measurements can help determine
$c_{\mathrm{ion},0}$, $\mu_{\mathrm{ion}}$, $D_{\mathrm{ion}}$, boundary
conditions, and the link between ion motion and defect growth. These
measurements are needed because ordinary J--V curves alone do not uniquely
determine the ion state.
Future DIT simulations should pair each ion-transport case with a control case
where the ions are held fixed, so that $I_{\mathrm{corr}}(t)$ and the extracted
$N_0$ are defined by the ion-related transient signal rather than by the total
measured current.

The fourth step is device-stack-specific calibration. The model should be
calibrated separately for device stacks, contacts, and barrier layers because
the same degradation pathway can produce different electrical signatures in
different architectures.

The fifth step is completion of the spatial and transient result set. The
selected spatial results now support band diagrams with carrier energy levels,
the correct HTL-to-ETL plotting direction, and selected ion and trap
line-profile panels from the profile grid. Final validation should still
include explicit short-circuit, maximum-power, and open-circuit band diagrams,
a strict COMSOL ion-conservation check, and fixed-ion-control DIT charge
extraction. The quantitative DIT result remains a future validation step until
a fixed-ion or no-ion control current is available for baseline subtraction.

The final step is a COMSOL-calibrated surrogate workbench. The current
6 $\times$ 6 light-temperature grid supports a first forward surrogate
baseline. That baseline demonstrates grouped COMSOL-to-COMSOL interpolation
for PCE, normalized PCE, PCE retention, $J_{sc}$, and J--V current density, but
it is trained only on the present COMSOL grid and is not an experimental
predictor. A later tool would need uncertainty estimates, active-learning run
selection, inverse diagnosis from measured J--V history, and matched
experimental calibration. The surrogate would speed up parameter exploration
and help infer hidden device states from measured J--V history, but it would
not replace the COMSOL drift--diffusion model as the physics reference.

\section{Concluding Perspective}

This thesis develops a COMSOL-based degradation model for perovskite solar
cells. The model links electronic transport, ion movement, defect growth,
contact resistance, leakage, voltage protocol, optical generation, DIT pulse
analysis, parameter sweeps, and accelerated J--V aging studies in one workflow.
It shows that several hidden mechanisms can produce distinct but overlapping
J--V signatures. Experimental calibration remains the main unfinished step.
Future work should focus on improved calibration, self-healing and
light-soaking terms, ion-sensitive validation, device-stack-specific
parameters, and broader COMSOL-calibrated surrogate use after the data set is
physically grounded.
